\section{From \emph{Summorum Pontificum} to the End of PCED}

Benedict XVI issued \emph{Summorum Pontificum} on 7 July 2007, while the Institute was still of diocesan right. The act called the 1962 Missal an extraordinary expression of the same Roman \emph{lex orandi}, broadened use by Latin priests, provided routes for stable groups, recognized community use in institutes and societies, and continued the Pontifical Commission \emph{Ecclesia Dei}. The following year's ICKSP decree used this vocabulary and located its solemn older worship within that legal environment.

The act did not found or erect the Institute. The ICKSP dated diocesan erection to 1990 and received pontifical-right status from N. 181/2008. Nor did \emph{Summorum Pontificum} make diocesan reception irrelevant. Article 3 referred community use to superiors and proper law, while public apostolates and personal parishes remained embedded in dioceses. The broader regime made requests and new works easier to normalize, but it did not create personnel or property.

\subsection{Normalization and a new generation}

Under the 2007 framework, faithful could encounter the older Mass without direct involvement in the 1970s and 1980s conflicts. ICKSP churches could function as stable liturgical, devotional, educational, and cultural centers rather than emergency accommodations. The society's seminary and canonial presentation offered bishops a formed community able to staff such works.

The checked set contains no survey of why candidates entered or why faithful attended. It supports an institutional inference: as time passed, the founding story became inherited identity rather than each member's personal experience. Gricigliano, common patrons, and the public history transmitted that identity across new cohorts.

\subsection{The Commission's suppression in 2019}

Francis decommissioned PCED on 17 January 2019 and transferred its tasks to the Congregation for the Doctrine of the Faith. The act noted that institutes formerly under the Commission had achieved stability of number and life and judged that remaining questions were principally doctrinal. The original Roman office that had erected and visited the ICKSP therefore disappeared.

This was an administrative and supervisory transition, not suppression of the Institute. The society, constitutions, houses, and members did not cease because their founding dicastery ended. But institutional maturity also meant integration into changing curial structures. In 2021 competence for former-PCED societies moved again, this time to the dicastery responsible for institutes of consecrated life and societies of apostolic life.

\statusframe{Benedict XVI, \emph{Summorum Pontificum}, 7 July 2007; Francis, motu proprio on PCED, 17 January 2019.}{The first act broadened the general older-liturgical regime before pontifical erection; the second ended the Commission after former-PCED institutes had achieved institutional stability.}{Neither act alone supplies an ICKSP-specific faculty beyond its text, and neither makes general government independent of bishops in public apostolate.}
